\documentclass[10pt,letterpaper]{article}
\usepackage{cornell-notes}

\title{Clause 03.47: Reserved Function}
\author{}
\date{}
\setDocTitle{Clause 03.47: Reserved Function}
\setDocSubtitle{C++ 2024}
\setDocOwner{C++ 2024 Cornell Notes}
\setDocDate{\today}
\setCornellCollection{C++ 2024 Cornell Notes}
\setCornellUnitType{Clause}
\setCornellUnitNumber{03.47}
\setCornellUnitTitle{Reserved Function}
\hypersetup{pdftitle={Clause 03.47: Reserved Function},pdfsubject={Cornell notes},pdfauthor={}}

\begin{document}
\maketitle
\begin{CornellOverviewBox}{Overview}
\textbf{Classification:} Definition\\[2pt]
\textbf{Learning objective:} Explain the precise meaning and rule context of reserved function.
\end{CornellOverviewBox}\begin{CornellSummaryBox}{Summary}
{\sffamily\bfseries\color{Accent} Summary}\par\vspace{3pt}Section 3.47 focuses on reserved function. Study the exact meaning, scope qualifier, and cross-references of the defined term; nearby terms can carry materially different semantics. Apply the section together with its cited rules so that terminology and constraints are interpreted in context.
\end{CornellSummaryBox}

\end{document}
